\begin{abstract}
Constraint solving on complex \SMT{} constraints such as nonlinear arithmetic, modular operations, and mixed logics remains challenging, as combinatorial blow-ups induce vast case splits and weak propagation under tight budgets. To address this challenge, we adopt a different approach: by concretizing a subset of the variables, we significantly reduce the complexity of the constraints. To efficiently identify the subset and its values, we present \tool, a solver-agnostic pre-solver that learns a policy for semantics-aware concretization. A \RL{} agent selects high-impact variables; a \LLM{} proposes semantics-guided values; the solver verifies each proposal in the loop. This decoupled design reduces branching and inter-variable coupling before the main solve and improves efficiency without modifying solver internals. 
Evaluations on SMTimer and \SMTCOMP{}, covering \QFLIA{} and \QFNIA{} with Z3, CVC5, MathSAT, and BVParti, show consistent gains. On SMTimer with Z3, the number of successfully solved instances rises from 72 to 94, a 30.6\% increase, and the average time per solved instance drops from 561 s to 214 s, a 2.6$\times$ speedup. On \QFNIA{}, with the same 1200s budget, \tool enables Z3 to solve 324 instances that previously timed out with baseline Z3. These results indicate that learned, policy-driven, semantics-aware concretization can turn otherwise intractable formulas into solvable ones while preserving solver independence.


\end{abstract}